% !TeX root = ../../Book3.tex
% 纲要标题：### 4.1 关联素理想
% 纲要位置：工作记录/计划纲要.md，第 2007 行。
% 纲要细目（写作范围）：
% * 元素零化子的素性；循环子模与素商模的嵌入、短正合列的非对称性
% * 关联素理想的存在和有限性；非典范素理想滤过、局部化中的统一清分母
% * 零因子作为关联素理想的并；支集极小点与孤立、嵌入关联素理想

\section{关联素理想}
\label{b3:sec:0401}

同一个元素可以被许多环元素零化，但它的全部零化子组成一个确定的理想。对交换含幺环 \(R\) 的模 \(M\) 及元素 \(m\in M\)，映射 \(a\mapsto am\) 给出 \(R\) 模同构 \(Rm\cong R/\operatorname{Ann}_R(m)\)。当这个零化子是素理想 \(\mathfrak p\) 时，模中便含有一个与 \(R/\mathfrak p\) 同构的循环子模。这给出了从模中提取素理想的方法。在 Noether 环上，每个非零模都有这样的素理想，有限生成模的关联素理想还只有有限个。前一个结论不要求模有限生成，而这些假设也不是每个具体模存在关联素理想的必要条件。

\subsection{元素零化子的素性}
\label{b3:subsec:0401:01}

\begin{definition}\label{b3:def:associated-prime}
设 \(R\) 为交换含幺环，\(M\) 为 \(R\) 模，\(m\in M\)。元素 \(m\) 的零化子为
\[
\operatorname{Ann}_R(m)=\{\,a\in R\mid am=0\,\}.
\]
若 \(R\) 的素理想 \(\mathfrak p\) 等于 \(M\) 中某个非零元素的零化子，则称 \(\mathfrak p\) 为 \(M\) 的\term[guanlian-sulixiang]{关联素理想}[associated prime]。所有这样的素理想组成集合 \(\operatorname{Ass}_R(M)\)，并约定 \(\operatorname{Ass}_R(0)=\varnothing\)。
\symidx{\(\operatorname{Ass}_R(M)\)：模的关联素理想集合}
\end{definition}

因此 \(\mathfrak p\in\operatorname{Ass}_R(M)\) 等价于存在 \(R\) 模单射 \(R/\mathfrak p\hookrightarrow M\)。定义说的是某个元素的零化子，不是整个模的零化子必须为素理想。

\begin{lemma}\label{b3:lem:maximal-element-annihilator-prime}
设 \(R\) 为交换含幺环，\(M\) 为 \(R\) 模，\(0\ne m\in M\)。若 \(\operatorname{Ann}_R(m)\) 在集合 \(\{\operatorname{Ann}_R(x):0\ne x\in M\}\) 中按包含关系极大，则它是 \(R\) 的素理想。
\end{lemma}
\begin{proof}
记 \(I=\operatorname{Ann}_R(m)\)。因 \(m\ne0\)，\(I\ne R\)。若 \(ab\in I\) 且 \(b\notin I\)，则 \(bm\ne0\)。乘以 \(b\) 后，原来零化 \(m\) 的元素仍零化 \(bm\)，所以 \(I\subseteq\operatorname{Ann}_R(bm)\)。由于 \(bm\) 仍是非零元素，极大性不允许这个零化子严格增大，只能有 \(\operatorname{Ann}_R(bm)=I\)。此时 \(a(bm)=0\) 说明 \(a\) 原来就属于 \(I\)，这正是素性。
\end{proof}

这里的“极大”指元素零化子这一族中的极大元，并不表示 \(I\) 是环的极大理想。例如整环 \(R\) 作为自身上的模时，每个非零元素的零化子都是 \((0)\)，于是 \(\operatorname{Ass}_R(R)=\bigl\{(0)\bigr\}\)，即使 \(R\) 不是域也成立。

\begin{proposition}\label{b3:prop:associated-exact-sequence}
设 \(R\) 为交换含幺环，\(0\longrightarrow N\longrightarrow M\longrightarrow P\longrightarrow0\) 为 \(R\) 模的短正合列。则
\[
\operatorname{Ass}_R(N)\subseteq\operatorname{Ass}_R(M)
 \subseteq\operatorname{Ass}_R(N)\cup\operatorname{Ass}_R(P).
\]
特别地，有限直和的关联素理想集合是各因子对应集合的并。对一般短正合列，却不能保证 \(\operatorname{Ass}_R(P)\subseteq\operatorname{Ass}_R(M)\)。
\end{proposition}
\begin{proof}
子模中的元素在 \(N\) 与 \(M\) 内有同一个零化子，得第一个包含。给定 \(\mathfrak p\in\operatorname{Ass}_R(M)\)，在 \(M\) 内取 \(L\cong R/\mathfrak p\)。\(L\) 中每个非零元素的零化子都是 \(\mathfrak p\)，因为 \(R/\mathfrak p\) 是整环。如果 \(L\cap N\ne0\)，取交内非零元素即得 \(\mathfrak p\in\operatorname{Ass}_R(N)\)；如果交为零，\(L\) 单射到 \(P\)，故 \(\mathfrak p\in\operatorname{Ass}_R(P)\)。直和结论由分裂正合列和两因子各自的嵌入得到，再对因子数归纳。

商模可以产生原模没有的关联素理想。例如，设 \(k\) 为域，\(R=k[t]\)，考虑短正合列
\[
0\longrightarrow tR\longrightarrow R\longrightarrow R/(t)\longrightarrow0.
\]
环 \(R\) 是整环，故 \(\operatorname{Ass}_R(R)=\{(0)\}\)；而 \(R/(t)\cong k\) 的每个非零元素的零化子均为 \((t)\)，所以 \(\operatorname{Ass}_R(R/(t))=\{(t)\}\)。因此商模的关联素理想不必属于中间模的关联素理想集合。
\end{proof}

\subsection{关联素理想的存在性与有限性}
\label{b3:subsec:0401:02}

\begin{theorem}\label{b3:thm:associated-existence}
若 \(R\) 是交换含幺 Noether 环，\(M\ne0\) 是 \(R\) 模，则 \(\operatorname{Ass}_R(M)\ne\varnothing\)。这里不要求 \(M\) 有限生成。
\end{theorem}
\begin{proof}
所有非零元素零化子构成非空理想族。Noether 升链条件保证其中有极大元，再用\cref{b3:lem:maximal-element-annihilator-prime}。
\end{proof}

Noether 环上的有限生成模可以逐层拆成素商模，再由短正合列的包含关系控制它的关联素理想。每一步的素商子模需要作出选择，所得滤过通常不唯一，其中出现的素理想只给关联素理想集合一个上界。

\begin{theorem}\label{b3:thm:associated-finite}
设 \(R\) 是交换含幺 Noether 环，\(M\) 是有限生成 \(R\) 模。存在有限子模链
\[
0=M_0\subsetneq M_1\subsetneq\cdots\subsetneq M_r=M,
\qquad M_i/M_{i-1}\cong R/\mathfrak p_i,
\]
其中各 \(\mathfrak p_i\) 为素理想；这种链称为\term[sulixiang-lvguo]{素理想滤过}[prime filtration]。并且
\[
\operatorname{Ass}_R(M)\subseteq\{\mathfrak p_1,\ldots,\mathfrak p_r\},
\]
所以关联素理想只有有限个。\(M=0\) 时采用空滤过。
\end{theorem}
\begin{proof}
若已构造 \(M_i\ne M\)，对非零商模 \(M/M_i\) 应用\cref{b3:thm:associated-existence}，取其中同构于 \(R/\mathfrak p_{i+1}\) 的子模，并令 \(M_{i+1}\) 为其原像。由\cref{b3:prop:noether-finite-modules}，有限生成模 \(M\) 是 Noether 模，故这个严格上升过程必在有限步停止。对短正合列
\(0\rightarrow M_{i-1}\rightarrow M_i\rightarrow R/\mathfrak p_i\rightarrow0\)
逐次应用\cref{b3:prop:associated-exact-sequence}，再用 \(\operatorname{Ass}_R(R/\mathfrak p_i)=\{\mathfrak p_i\}\)，得到所需包含。
\end{proof}

例如 \(\mathbb Z\supset2\mathbb Z\supset0\) 的两个因子为 \(\mathbb Z/2\mathbb Z\) 与 \(\mathbb Z\)，但 \(\operatorname{Ass}_{\mathbb Z}(\mathbb Z)=\bigl\{(0)\bigr\}\)，滤过中的 \((2)\) 是额外出现的素理想。此外，有限生成性不能删除：\(\bigoplus_{p\;\text{为素数}}\mathbb Z/p\mathbb Z\) 的关联素理想包括所有 \((p)\)，因而是无限集。

\begin{proposition}\label{b3:prop:associated-localization}
设 \(R\) 为交换含幺 Noether 环，\(S\subseteq R\) 为含一的乘法集，\(M\) 为任意 \(R\) 模。则
\[
\operatorname{Ass}_{S^{-1}R}(S^{-1}M)
=\bigl\{\,S^{-1}\mathfrak p\mid \mathfrak p\in\operatorname{Ass}_R(M),\quad
                       \mathfrak p\cap S=\varnothing\,\bigr\}.
\]
\end{proposition}
\begin{proof}
若 \(\operatorname{Ann}_R(m)=\mathfrak p\) 且 \(\mathfrak p\cap S=\varnothing\)，则 \(m/1\ne0\)。当 \((a/s)(m/1)=0\) 时，有 \(t\in S\) 使 \(tam=0\)；由 \(ta\in\mathfrak p\)、\(t\notin\mathfrak p\) 得 \(a\in\mathfrak p\)。故其零化子恰为 \(S^{-1}\mathfrak p\)。

反之，设 \(\mathfrak q\) 是非零元素 \(m/s\) 的素零化子，令 \(\mathfrak p\) 为 \(\mathfrak q\) 在 \(R\) 中的收缩。由\cref{b3:thm:spec-localization}，\(\mathfrak p\) 是避开 \(S\) 的素理想，且 \(\mathfrak q=S^{-1}\mathfrak p\)。以单位 \(s/1\) 乘此元素后可改用 \(m/1\)。要在原模中找到零化子恰为 \(\mathfrak p\) 的元素，必须把局部成立的零化关系同时变成原模中的等式；这一步需要共同的分母。由于 \(R\) 是 Noether 环，理想 \(\mathfrak p\) 有有限生成集，写成 \(\mathfrak p=(a_1,\ldots,a_l)\)。因 \((a_i m)/1=0\)，可选 \(s_i\in S\) 使 \(s_i a_i m=0\)。令 \(t=s_1\cdots s_l\)，此时 \(\mathfrak p\subseteq\operatorname{Ann}_R(tm)\)；当 \(\mathfrak p=(0)\) 时采用空生成集并令 \(t=1\)。若 \(a(tm)=0\)，由于 \(t/1\) 可逆，局部化后 \(a/1\) 便零化 \(m/1\)，所以 \(a\in\mathfrak p\)。于是 \(\operatorname{Ann}_R(tm)=\mathfrak p\)，且 \(tm\ne0\)，否则 \(m/1=0\)。这就找到了原模中的关联元素，而所用的有限生成性来自理想 \(\mathfrak p\)，并未要求 \(M\) 有限生成。
\end{proof}

例如，设 \(k\) 为域，\(R=k[x,y]\)、\(M=R/(x)\)、\(\mathfrak m=(x,y)\)。因 \(R/(x)\cong k[y]\) 是整环，每个非零模元素的零化子都是 \((x)\)，故 \(\operatorname{Ass}_R(M)=\{(x)\}\)。局部化后有
\[
\begin{aligned}
M_{\mathfrak m}&\cong R_{\mathfrak m}/(x)R_{\mathfrak m}
\cong k[y]_{(y)},\\
\operatorname{Ass}_{R_{\mathfrak m}}(M_{\mathfrak m})
&=\{(x)R_{\mathfrak m}\}.
\end{aligned}
\]
这个唯一的关联素理想严格小于局部环的极大理想 \(\mathfrak mR_{\mathfrak m}\)：元素 \(y/1\) 属于后者，在商环 \(k[y]_{(y)}\) 中的像却非零，所以不属于前者。

\begin{proposition}\label{b3:prop:minimal-primes-associated}
设 \(R\) 为交换含幺 Noether 环，\(M\) 为非零有限生成 \(R\) 模。则 \(\operatorname{Supp}_R(M)\) 的每个素理想都包含某个关联素理想，且
\[
\min\operatorname{Supp}_R(M)=\min\operatorname{Ass}_R(M).
\]
尤其，对任意真理想 \(I\)，包含 \(I\) 的极小素理想都是 \(R/I\) 的关联素理想。
\end{proposition}
\begin{proof}
若 \(\mathfrak p\in\operatorname{Ass}_R(M)\)，上述局部化命题保证 \(M_{\mathfrak p}\ne0\)，故 \(\operatorname{Ass}_R(M)\subseteq\operatorname{Supp}_R(M)\)。给定 \(\mathfrak q\in\operatorname{Supp}_R(M)\)，\(R_{\mathfrak q}\) 是 Noether 环，非零模 \(M_{\mathfrak q}\) 有关联素理想；由\cref{b3:prop:associated-localization}，它来自某个 \(\mathfrak p\in\operatorname{Ass}_R(M)\)，且 \(\mathfrak p\subseteq\mathfrak q\)。因此两集合的极小元一致。对有限生成模，\cref{b3:thm:support-finite}给出 \(\operatorname{Supp}_R(M)=V\bigl(\operatorname{Ann}_R(M)\bigr)\)。取 \(M=R/I\) 就得到最后一句。
\end{proof}

设 \(k\) 为域，取 \(R=k[x,y]\)、\(M=R/(x^2,xy)\)。由 \(\sqrt{(x^2,xy)}=(x)\)，其支集为 \(V(x)\)，极小元只有 \((x)\)，所以 \((x)\) 也是极小关联素理想。另一方面，非零元素 \(\bar x\) 的零化子为 \((x,y)\)，具体计算见\subsecref{b3:subsec:0404:02}。因此 \((x,y)\) 也是关联素理想，却严格包含 \((x)\)。

\subsection{零因子的关联素理想描述}
\label{b3:subsec:0401:03}

设 \(R\) 为交换含幺环，\(M\) 为 \(R\) 模。若存在 \(0\ne m\in M\) 使 \(am=0\)，则称 \(a\in R\) 是 \(M\) 上的\term[mo-shang-lingyinzi]{零因子}[zero divisor on a module]；否则称它在 \(M\) 上是正则的。记这些零因子的集合为 \(Z_R(M)\)。这里允许 \(a=0\)：对非零模，零当然属于这个集合；对零模，集合为空。

\ProseParagraphBreak
标量 \(a\) 的乘法映射 \(M\to M\)、\(m\mapsto am\) 单射，当且仅当 \(a\notin Z_R(M)\)。因此零因子集合决定哪些标量在模上的作用失去单射性；在 Noether 环上，关联素理想恰好给出这个集合的描述。

\begin{theorem}\label{b3:thm:associated-zero-divisors}
设 \(R\) 为交换含幺 Noether 环，\(M\) 为 \(R\) 模。则
\[
Z_R(M)=\bigcup_{\mathfrak p\in\operatorname{Ass}_R(M)}\mathfrak p.
\]
若 \(M\) 有限生成，这就是有限个素理想的并。
\end{theorem}
\begin{proof}
关联素理想中的元素会零化定义它的非零元素，因此右端包含于左端。反向设 \(am=0\)、\(m\ne0\)。包含 \(\operatorname{Ann}_R(m)\) 的非零元素零化子构成非空理想族；由于 \(R\) 是 Noether 环，其中存在极大元 \(\operatorname{Ann}_R(n)\)。它在全部非零元素零化子中也是极大的，因任何更大零化子仍包含 \(\operatorname{Ann}_R(m)\)。\cref{b3:lem:maximal-element-annihilator-prime}说明它是关联素理想，而 \(a\) 在其中。
\end{proof}

例如，设 \(k\) 为域，在 \(R=k[x,y]/(xy)\) 中，\(\bar x\) 与 \(\bar y\) 的零化子分别为 \((\bar y)\) 和 \((\bar x)\)。每个元素都可唯一写为常数加一个无常数项的 \(x\) 多项式，再加一个无常数项的 \(y\) 多项式。限制到两轴给出单射 \(R\to k[x]\times k[y]\)；不在 \((\bar x)\cup(\bar y)\) 中的元素在两个坐标上都非零，故乘积为零时，两个整环中的消去律迫使另一乘子为零。而 \((\bar x)\)、\((\bar y)\) 中的元素分别被 \(\bar y\)、\(\bar x\) 零化，所以零因子恰为这两个素理想的并。\(\bar x+\bar y\) 不在其中，因而不是零因子，尽管它的两个加项都是零因子。这说明零因子集合不一定对加法封闭，因而不一定是理想。

\ProseParagraphBreak
更一般地，设 \(R\) 为交换含幺 Noether 环，若 \(R\) 的理想 \(J\) 的每个元素都是有限生成 \(R\) 模 \(M\) 上的零因子，由\cref{b3:thm:associated-finite}和\cref{b3:thm:associated-zero-divisors}及\cref{b3:lem:finite-prime-avoidance}的有限素理想回避性质，\(J\) 必包含于某一个关联素理想。这把“逐个检查乘法是否单射”的问题转成了有限个素理想的包含问题。以上存在性、局部化和滤过方法可对照 \cite[\S6, Theorems 6.1--6.5]{Matsumura1989CommutativeRingTheory}。

\ProseParagraphBreak
设 \(R\) 为交换含幺 Noether 环，\(M\) 为有限生成 \(R\) 模。按包含关系极小的关联素理想称为\term[guli-guanlian-sulixiang]{孤立关联素理想}[isolated associated prime]，其余称为\term[qianru-guanlian-sulixiang]{嵌入关联素理想}[embedded associated prime]。这里的“孤立”指素理想包含关系中的极小性，不要求它是 Zariski 拓扑中的孤立点，即不要求单点集合为开集。由\cref{b3:prop:minimal-primes-associated}，前者就是支集的极小元，后者则严格包含某个极小关联素理想；两类都由 \(M\) 本身确定。\cref{b3:tab:support-associated-primes}汇集了这些关系。

\begin{table}[htbp]
\centering
\caption[支集与关联素理想]{支集与关联素理想。设 \(R\) 为交换含幺 Noether 环，\(M\) 为有限生成 \(R\) 模；极小元均按素理想的包含关系理解。}
\label{b3:tab:support-associated-primes}
\begin{tabularx}{\textwidth}{@{}>{\raggedright\arraybackslash}p{0.25\textwidth}>{\raggedright\arraybackslash}X@{}}
\toprule
集合 & 与支集及零化子的关系 \\
\midrule
支集 & \(\operatorname{Supp}_R(M)=V\bigl(\operatorname{Ann}_R(M)\bigr)\)；每个支集中的素理想都包含某个关联素理想 \\
关联素理想 & \(\operatorname{Ass}_R(M)\subseteq\operatorname{Supp}_R(M)\)，且该集合有限 \\
孤立关联素理想 & \(\min\operatorname{Ass}_R(M)=\min\operatorname{Supp}_R(M)\)，即支集的极小点 \\
嵌入关联素理想 & \(\operatorname{Ass}_R(M)\setminus\min\operatorname{Supp}_R(M)\)，即支集中非极小的关联点 \\
\bottomrule
\end{tabularx}
\end{table}
